%==========================================================
%  Introduction.tex  —  Introduction générale
%==========================================================
\chapter*{Introduction générale}
\addcontentsline{toc}{chapter}{Introduction générale}
\markboth{Introduction générale}{}

\thispagestyle{fancy}

\definecolor{tekupBlue}{RGB}{0,70,127}
\definecolor{lightBlue}{RGB}{230,240,250}

%----------------------------------------------------------
% Paragraphe 1 — Contexte général du CRM et IA
%----------------------------------------------------------
La transformation numérique des centres d'appels constitue aujourd'hui un enjeu
majeur pour les entreprises cherchant à améliorer la qualité de leur relation
client. Dans un environnement où le volume d'appels ne cesse de croître et où
les attentes des clients en matière de réactivité et de personnalisation
s'intensifient, les centres de contact traditionnels doivent évoluer vers des
plateformes intelligentes, capables d'optimiser chaque interaction.
L'intelligence artificielle, et plus particulièrement les modèles de langage
(LLM) et la reconnaissance vocale (ASR), ouvrent des perspectives inédites
pour l'analyse des appels, l'évaluation de la qualité et l'aide à la
décision.

%----------------------------------------------------------
% Paragraphe 2 — Contexte entreprise et objectif du projet
%----------------------------------------------------------
C'est dans cette perspective que s'inscrit ce projet de fin d'études, réalisé
au sein de \textbf{3LM Solution}, entreprise de services numériques tunisienne
spécialisée dans les solutions cloud-natives et la transformation digitale.
3LM Solution a exprimé le besoin d'une plateforme CRM (Customer Relationship
Management) moderne, intégrant des capacités d'intelligence artificielle pour
automatiser et enrichir le suivi de la relation client dans un contexte de
centre d'appels. L'objectif de ce projet est donc la \textbf{conception et le
développement d'un CRM intelligent} couvrant la gestion des leads, des appels,
des rendez-vous, des campagnes, de la qualité, de la paie et de la présence,
le tout enrichi par un moteur de scoring IA.

%----------------------------------------------------------
% Paragraphe 3 — Réalisation technique et méthodologie
%----------------------------------------------------------
J'ai conçu et développé seul cette plateforme, en m'appuyant sur une
architecture moderne avec un backend en \textbf{.NET 8} et \textbf{ASP.NET
Core} utilisant le pattern Repository et Unit of Work, et un frontend en
\textbf{React 18} et \textbf{TypeScript} avec l'écosystème \textbf{Radix UI}
et \textbf{TailwindCSS}. Le moteur d'IA s'appuie sur \textbf{Ollama} avec le
modèle \textbf{gemma3:4b} pour l'analyse des appels, \textbf{Whisper} pour la
transcription automatique et \textbf{DistilBERT} pour l'analyse de sentiment.
La base de données est assurée par \textbf{PostgreSQL 16}. Le développement a
été mené selon une démarche \textbf{Agile Scrum}, en trois releases
successives totalisant six sprints : la \textbf{Release 1} a posé les
fondations (authentification, gestion des utilisateurs, leads et rendez-vous)
; la \textbf{Release 2} a ajouté le cœur IA (scoring, transcription, analyse
de qualité) ; et la \textbf{Release 3} a consolidé la solution (présence,
paie, rapports, déploiement continu).

%----------------------------------------------------------
% Paragraphe 4 — Structure du rapport
%----------------------------------------------------------
Ce rapport est organisé en six chapitres~: le \textbf{Chapitre 1} présente
le cadre du projet et l'organisme d'accueil~; le \textbf{Chapitre 2} détaille
l'analyse des besoins et l'architecture globale retenue~; le \textbf{Chapitre
3} décrit l'environnement de travail et les technologies mobilisées~; le
\textbf{Chapitre 4} présente la Release 1 (authentification, utilisateurs,
leads et rendez-vous)~; le \textbf{Chapitre 5} expose la Release 2 (scoring IA,
transcription et qualité)~; et le \textbf{Chapitre 6} finalise avec la Release 3
(présence, paie, analytics et déploiement). Une \textbf{conclusion générale}
clôt ce rapport en synthétisant les apports du projet et ses perspectives
d'évolution.
